%% kapa: exam=1 round=37 year=2026 period=1 form=A subject=civil_law subject_ko=민법 seq=12 number=12
%% source: Q-Net 감정평가사 기출문제 (cst00309) artlSeq=5258924 file=raw/1008/5258924/2026년 제37회 감정평가사 1차 1교시.pdf
%% conversion: kordoc 4.24.51 -> tools/gen1.py (auto); review=proofread-sonnet; flags=none
\begin{question}{12}
착오에 의한 의사표시에 관한 설명으로 옳지 않은 것은? (다툼이 있으면 판례에 따름)
\begin{choices}
  \item 소송행위에는 착오에 의한 의사표시의 취소가 인정되지 않는다.
  \item 대리인에 의해 계약이 체결된 경우, 착오의 유무는 대리인을 기준으로 판단하여야 한다.
  \item 매도인이 매수인의 채무불이행을 이유로 매매계약을 적법하게 해제한 경우, 착오를 이유로 한 매수인의 취소권은 소멸한다.
  \item 착오로 인하여 표의자가 경제적 불이익을 입지 않는 경우에는 특별한 사정이 없는 한 중요부분의 착오라고 할 수 없다.
  \item 매매계약 내용의 중요부분에 착오가 있는 경우, 매수인은 매도인의 하자담보책임이 성립하는지와 상관없이 착오를 이유로 그 매매계약을 취소할 수 있다.
\end{choices}
\end{question}
\answer{3}
